\documentclass[11pt, tikz, border=12pt]{standalone}
\usepackage[fontset=fandol]{ctex}
\usetikzlibrary{positioning, shadows, arrows.meta, backgrounds, fit, calc}

\definecolor{tcyan}{RGB}{72, 122, 138}
\definecolor{tyellow}{RGB}{178, 132, 76}
\definecolor{tgreen}{RGB}{88, 138, 118}
\definecolor{tpink}{RGB}{156, 96, 116}

\begin{document}
\begin{tikzpicture}[
    font=\sffamily,
    box/.style={
        draw=#1,
        fill=#1!6,
        thick,
        rounded corners=3pt,
        text width=4.3cm,
        minimum height=1.55cm,
        align=center,
        inner sep=6pt,
        font=\footnotesize
    },
    header/.style={
        fill=#1,
        text=white,
        font=\bfseries\small,
        minimum width=4.55cm,
        inner sep=5pt,
        rounded corners=2pt,
        align=center
    },
    labelband/.style={
        fill=#1!10,
        draw=#1!85,
        dashed,
        thick,
        rounded corners=3pt,
        align=center,
        font=\footnotesize\bfseries,
        inner sep=6pt,
        minimum height=1.05cm
    },
    arrow/.style={-{Latex[scale=1.0]}, thick, draw=black!55},
    flowarrow/.style={-{Latex[scale=1.15]}, line width=1.3pt, draw=black!45}
]

    % ==================== 1. 第一列：CPT ====================
    \node[header=tcyan] (cpt-head) at (0, 0) {阶段一：CPT 知识注入};
    \node[box=tcyan, below=0.35cm of cpt-head] (cpt-in) {\textbf{输入层}\\通用基座（Qwen2.5-7B）\\七类电力语料 + 通用语料};
    \node[box=tcyan, below=0.45cm of cpt-in] (cpt-proc) {\textbf{核心处理}\\PPL 双阈值筛选（$P_{10}$--$P_{80}$）\\1:1 混合因果语言建模};
    \node[box=tcyan, below=0.45cm of cpt-proc] (cpt-out) {\textbf{输出层}\\CPT 电力基础模型\\（具备领域事实先验）};

    % ==================== 2. 第二列：SFT ====================
    \node[header=tyellow, right=1.15cm of cpt-head] (sft-head) {阶段二：SFT 任务微调};
    \node[box=tyellow, below=0.35cm of sft-head] (sft-in) {\textbf{输入层}\\CPT 输出基座模型\\约 6.8 万条电力指令数据};
    \node[box=tyellow, below=0.45cm of sft-in] (sft-proc) {\textbf{核心处理}\\六模块场景原生合成管线\\两阶段微调 + DFT 奖励校正};
    \node[box=tyellow, below=0.45cm of sft-proc] (sft-out) {\textbf{输出层}\\SFT 多任务模型\\（保留七类原生风格）};

    % ==================== 3. 第三列：RLHF ====================
    \node[header=tgreen, right=1.15cm of sft-head] (rlhf-head) {阶段三：RLHF 风格对齐};
    \node[box=tgreen, below=0.35cm of rlhf-head] (rlhf-in) {\textbf{输入层}\\SFT 输出基座模型\\约 1.5 万条风格对齐样本};
    \node[box=tgreen, below=0.45cm of rlhf-in] (rlhf-proc) {\textbf{核心处理}\\三维强度条件化 Prefix\\CORAL-7B 判别 + GRPO};
    \node[box=tgreen, below=0.45cm of rlhf-proc] (rlhf-out) {\textbf{输出层}\\RLHF 风格可控模型\\（显式调节三维风格强度）};

    % ==================== 4. 第四列：DPO ====================
    \node[header=tpink, right=1.15cm of rlhf-head] (dpo-head) {阶段四：DPO 安全对齐};
    \node[box=tpink, below=0.35cm of dpo-head] (dpo-in) {\textbf{输入层}\\RLHF 输出基座模型\\三类协同安全偏好数据};
    \node[box=tpink, below=0.45cm of dpo-in] (dpo-proc) {\textbf{核心处理}\\Clean / Recovery / Counter\\困难状态挖掘 + HS-Weighting};
    \node[box=tpink, below=0.45cm of dpo-proc] (dpo-out) {\textbf{输出层}\\最终部署模型\\（知识+任务+风格+安全）};

    % ==================== 5. 顶部标签数据流 ====================
    \coordinate (topY) at ([yshift=1.45cm]cpt-head.north);
    \path (cpt-head.west) -- (sft-head.east) coordinate[midway] (mid12);
    \node[labelband=tcyan, text width=9.8cm] (label7) at (mid12 |- topY) {七类电力数据标签贯穿（保留领域知识来源与原生业务场景）};
    \node[labelband=tgreen, text width=4.3cm] (label3) at (rlhf-head |- topY) {三类场景归纳标签\\（风格控制目标映射）};
    \node[labelband=tpink, text width=4.3cm] (labelsafe) at (dpo-head |- topY) {四类安全风险标签\\（深层安全合规边界）};

    \draw[arrow, draw=tcyan!90!black] (label7.east) -- node[above, font=\scriptsize] {归纳} (label3.west);
    \draw[arrow, dashed, draw=black!50] (label3.east) -- (labelsafe.west);

    \draw[arrow, dashed, draw=tcyan!85!black] (label7.south -| cpt-head.north) -- (cpt-head.north);
    \draw[arrow, dashed, draw=tcyan!85!black] (label7.south -| sft-head.north) -- (sft-head.north);
    \draw[arrow, dashed, draw=tgreen!85!black] (label3.south) -- (rlhf-head.north);
    \draw[arrow, dashed, draw=tpink!85!black] (labelsafe.south) -- (dpo-head.north);

    % ==================== 6. 纵向与横向连接线 ====================
    \foreach \col in {cpt, sft, rlhf, dpo} {
        \draw[arrow] (\col-head) -- (\col-in);
        \draw[arrow] (\col-in) -- (\col-proc);
        \draw[arrow] (\col-proc) -- (\col-out);
    }

    \draw[flowarrow] (cpt-proc.east) -- (sft-proc.west);
    \draw[flowarrow] (sft-proc.east) -- (rlhf-proc.west);
    \draw[flowarrow] (rlhf-proc.east) -- (dpo-proc.west);

    \draw[arrow, dashed, draw=tcyan!90!black] (cpt-out.east) -- node[above, font=\scriptsize] {初始化} (sft-out.west);
    \draw[arrow, dashed, draw=tyellow!90!black] (sft-out.east) -- node[above, font=\scriptsize] {初始化} (rlhf-out.west);
    \draw[arrow, dashed, draw=tgreen!90!black] (rlhf-out.east) -- node[above, font=\scriptsize] {初始化} (dpo-out.west);

    \begin{scope}[on background layer]
        \node[fit=(label7)(labelsafe)(cpt-head)(dpo-out), draw=black!12, fill=gray!3, rounded corners=6pt, inner sep=12pt] {};
    \end{scope}

\end{tikzpicture}
\end{document}
